\documentclass[11pt]{article}
% Draft of paper 1 (the bound). Reshaped from docs/dnn_bound/dnn_bound.tex (Nov 5 2026),
% whose branch-and-bound and nug30 material moves to the companion paper.
% Conventions in this draft:
%   \todo{...}   something missing or to be decided (red in the PDF)
%   % SOURCE:    where in the repository a number comes from
\usepackage[letterpaper,margin=1in]{geometry}
\usepackage[T1]{fontenc}
\usepackage{lmodern}
\usepackage{amsmath,amssymb,amsthm}
\usepackage{booktabs}
\usepackage{array}
\usepackage[labelfont=bf,labelsep=period,font=small]{caption}
\usepackage{xcolor}
\usepackage[hidelinks]{hyperref}

\newtheorem{proposition}{Proposition}
\newtheorem{lemma}[proposition]{Lemma}
\theoremstyle{remark}
\newtheorem{remark}[proposition]{Remark}

\newcommand{\R}{\mathbb{R}}
\newcommand{\Sym}{\mathcal{S}}
\newcommand{\tr}{\operatorname{tr}}
\newcommand{\vecop}{\operatorname{vec}}
\newcommand{\inner}[2]{\langle #1, #2\rangle}
\newcommand{\lmax}{\lambda_{\max}}
\newcommand{\Pspd}{\mathcal{P}_{+}}
\newcommand{\QPB}{\mathrm{QPB}}
\newcommand{\DNN}{\mathrm{DNN}}
\newcommand{\DNNp}{\mathrm{DNN}^{+}}
\newcommand{\GL}{\mathrm{GL}}
\newcommand{\LAP}{\mathrm{LAP}}
\newcommand{\todo}[1]{{\color{red}[#1]}}

\title{An extended doubly nonnegative bound\\ for the quadratic assignment problem}
\author{Nathan W. Brixius \and Jean-Pierre Goux}
\date{Draft, November 24, 2026}

\begin{document}
\maketitle

\begin{abstract}
We describe a lower bound for the quadratic assignment problem (QAP) that extends
the doubly nonnegative (DNN) semidefinite relaxation in two ways. First, the bound
is \emph{certified}: it is the dual value of an ADMM multiplier whose projection
onto the relaxation's face is negative semidefinite by construction, evaluated by
a Gilmore--Lawler computation on the lifted cost matrix, less explicit
floating-point allowances, so that it is a valid bound after any number of
iterations. Second, the relaxation is strengthened by \emph{regional} triangle
inequalities: Boolean-quadric inequalities on the indicators that a facility is
assigned to a given set of locations. For sets that decompose the distance matrix
(row and column prefixes of a rectangular grid, or half-planes and clusters for
near-planar distances) these inequalities act directly on the terms of the
objective, are linear in the DNN matrix, and are not implied by the DNN
relaxation. They are enforced inside the ADMM by a cutting-plane projection, at a
cost of a few per cent per iteration, and every intermediate value remains
certified. We call the result $\DNNp$. At the root it improves the best published
lower bounds for several large QAPLIB problems, for example tho40 from 226,522 to
230,453, tai50b from 431,245,612 to 453,552,676 and tai60b from 592,829,061 to
605,315,696, and it closes 96\% or more of the DNN gap on tai25b and tai30b.
\todo{Final figures after the root runs of Section~\ref{sec:results} are redone
under one stopping rule; add sko42, sko49, wil50.} We also describe the properties
that make the bound usable in branch and bound: validity at every iteration, warm
starts from a parent's state, and cuts that are inherited by a node's children.
\end{abstract}

\section{Introduction}

The quadratic assignment problem (QAP) in ``Koopmans--Beckmann'' form is to
\[
  \min_{\pi}\ \sum_{i=1}^n\sum_{k=1}^n a_{ik}\,b_{\pi(i),\pi(k)} + \sum_{i=1}^n c_{i,\pi(i)},
\]
where $a_{ik}$ is the flow between facilities $i$ and $k$, $b_{jl}$ is the
distance between locations $j$ and $l$, $c_{ij}$ is the fixed cost of assigning
facility $i$ to location $j$, and $\pi(i) = j$ if facility $i$ is assigned to
location $j$. In matrix form,
\begin{equation}\label{eq:qap}
  \begin{array}{rl}
  \mathrm{QAP}(A,B,C): & \min\ \tr(AXB + C)X^T\\
  & \ \text{s.t.}\ \ X\in\Pi,
  \end{array}
\end{equation}
where $\Pi$ is the set of $n\times n$ permutation matrices. QAPLIB \cite{qaplib}
has long served as the benchmark library for the QAP; for surveys see
\cite{burkardsurvey,loiola}.

The QAP is NP-hard, and in practice is one of the hardest discrete optimization
problems to solve to optimality. Exact methods are almost exclusively branch and
bound, and their success rests on the lower bound: on its quality, on its cost,
and on how well it can be reused from one node of the tree to the next. The
Gilmore--Lawler bound (GLB) \cite{gilmore,lawler} is cheap but weak for larger
problems. Bounds based on eigenvalues \cite{hrw}, on convex quadratic programming
\cite{ab01}, on the reformulation-linearization technique \cite{datenagi} and on
semidefinite programming (SDP) \cite{zkrw} are progressively stronger and more
expensive. Among the bounds that have been computed for problems with $n\ge30$,
the strongest are those from the doubly nonnegative (DNN) relaxation of Zhao,
Karisch, Rendl and Wolkowicz \cite{zkrw}: the SDP relaxation of the lifted matrix
$yy^T$, $y = (1,\vecop(X))$, restricted to its minimal face, with the ``gangster''
zeros and elementwise nonnegativity. First-order methods made this relaxation
computable at $n = 30$ to $60$: the ADMM of Oliveira, Wolkowicz and Xu \cite{owx},
the restricted Peaceman--Rachford method of Graham et al.\ \cite{graham}, and the
bisection--projection and Newton-bracketing methods for the Lagrangian-DNN
relaxation \cite{bbcpop,kkt}. Fujii et al.\ \cite{fujii} used the last of these
at every node of a branch-and-bound tree on large parallel computers. The best
published lower bounds for the large unsolved QAPLIB problems are all values of
this one relaxation, differing in the last digits with the accuracy of the solver.

This paper extends the DNN bound in two directions.
\begin{itemize}
\item[(i)] \emph{A certified bound} (Section~\ref{sec:cert}). First-order iterates
  are neither primal feasible nor dual optimal, and published DNN bounds are
  obtained by running a solver to a tolerance and evaluating a dual function at the
  final iterate. We give a bound that is valid for \emph{any} multiplier, with
  explicit allowances for floating-point error, and a choice of multiplier and of
  polyhedral relaxation that make it tight and cheap to evaluate at every few
  iterations. The dual function is that of \cite{graham}; the multiplier, the
  two-level assignment evaluation and the error analysis are new.
\item[(ii)] \emph{Regional cuts} (Section~\ref{sec:cuts}). For a set $S$ of
  locations, the indicators $z_i = \sum_{a\in S}x_{ia}$ that facility $i$ is placed
  in $S$ are binary, and their first and second moments are linear functions of
  the DNN matrix. The triangle inequalities of the Boolean quadric polytope
  \cite{padberg} on these moments are valid for the QAP and not implied by the DNN
  relaxation. When the sets $S$ are chosen so that the distance matrix is a
  nonnegative combination of their cut semimetrics, as it is exactly for the
  Manhattan distances of grid problems and approximately for the clustered
  distances of Taillard's ``b'' problems, the inequalities constrain exactly the
  quantities that appear in the objective. Triangle inequalities on the entries of
  the lifted matrix itself, by contrast, gain almost nothing.
\item[(iii)] \emph{Cuts inside the ADMM} (Section~\ref{sec:cpadmm}). Following the
  cutting-plane ADMM of de Meijer et al.\ \cite{dmswz} for graph partition, the
  inequalities are imposed in the polyhedral step of the ADMM by Dykstra's
  algorithm, which here reduces to updates on the small moment matrices of the
  regions. The multipliers of that projection extend the certificate of (i), so
  that the bound with cuts is also valid at every iteration.
\end{itemize}
We call the resulting bound $\DNNp$. Section~\ref{sec:bnb} describes the
properties that matter for its use in branch and bound, and
Section~\ref{sec:results} compares it with published bounds. The validation of the
bound is described in Appendix~\ref{sec:validation}.

The bound was developed for a branch-and-bound code descended from that of
\cite{ba01,abgl}. That algorithm, its branching rules, and the solution to
optimality of several QAPLIB problems with it, including problems for which the
values listed here as ``best known'' are thereby proved optimal, are the subject
of a companion paper \cite{companion}. Here we consider only the bound.

\paragraph{Notation.} We use $\tr(\cdot)$ to denote the trace of a matrix. If $A$
and $B$ are matrices then $\inner{A}{B} = \tr(AB^T)$, $A\otimes B$ is the Kronecker
product of $A$ and $B$, and $\vecop(A)$ is the vector formed by ``stacking'' the
columns of $A$. We use $e$ to denote a vector with each component equal to one,
$\Sym^N$ for the symmetric $N\times N$ matrices, and $\lmax(\cdot)$ for the largest
eigenvalue of a symmetric matrix. For convenience we let the name of an
optimization problem, like $\mathrm{QAP}(A,B,C)$, also refer to its optimal value.

\section{The DNN relaxation}\label{sec:dnn}

Consider a node of the branch-and-bound tree with $m$ facilities left to assign.
Its subproblem is again a QAP \eqref{eq:qap} of size $m$, with reduced matrices
$A, B \in \R^{m\times m}$, a linear term $C$ that collects the interactions with the
fixed assignments, and a constant $s$. For $X\in\Pi$ let $x = \vecop(X^T)$, ordered
so that $x_{ij}$ (facility $i$ at location $j$) has lifted index $1 + im + j$, and
let $y = (1, x)$ and $N = m^2+1$. Define the symmetric lifted cost $L\in\Sym^N$ by
\begin{equation}\label{eq:L}
  L_{(ij),(kl)} = \tfrac12(a_{ik}b_{jl} + a_{ki}b_{lj}), \qquad
  L_{0,(ij)} = L_{(ij),0} = \tfrac12 c_{ij}, \qquad L_{00} = 0.
\end{equation}
Then for every permutation, $\inner{L}{yy^T} = \sum a_{ik}b_{jl}x_{ij}x_{kl} +
\sum c_{ij}x_{ij}$, the subproblem objective without $s$, for any $A$ and $B$,
symmetric or not.

\paragraph{The face.} The constraints $Xe = X^Te = e$ say that $y$ lies in the null
space of a linear map, so every lifted permutation $Y = yy^T$ has its range in a
fixed subspace of dimension $(m-1)^2+1$. Let $V \in\R^{N\times K}$, $K = (m-1)^2+1$,
be an orthonormal basis of it. Zhao et al.\ \cite{zkrw} take
$V = \begin{pmatrix}1 & 0\\ e/m & V_m\otimes V_m\end{pmatrix}$, suitably normalized,
where the columns of $V_m\in\R^{m\times(m-1)}$ span $e^\perp$. Every lifted
permutation then satisfies $Y = VRV^T$ with $R \succeq 0$. The DNN relaxation of the
subproblem is
\begin{equation}\label{eq:dnn}
  \begin{array}{rl}
  \DNN(A,B,C) = s + \min & \inner{L}{Y}\\
  \text{s.t.} & Y = VRV^T,\ R\succeq 0,\ Y_{00}=1,\ \mathcal{G}(Y) = 0,\ 0 \le Y \le 1,
  \end{array}
\end{equation}
where $\mathcal{G}$ selects the ``gangster'' entries $Y_{(ij),(il)}$, $l\ne j$, and
$Y_{(ij),(kj)}$, $k\ne i$, which vanish for permutations since a facility has one
location and a location one facility. Oliveira, Wolkowicz and Xu \cite{owx} solve
\eqref{eq:dnn} by ADMM on the splitting $Y = VRV^T$, alternating
\begin{align}
  R &\leftarrow \Pspd\big(V^T(Y + Z/\beta)V\big), \label{eq:rstep}\\
  Y &\leftarrow \Pi_{\mathcal{Y}}\big(VRV^T - (L + Z)/\beta\big), \label{eq:ystep}\\
  Z &\leftarrow Z + \gamma\beta\,(Y - VRV^T), \label{eq:zstep}
\end{align}
where $\Pspd$ projects onto the positive semidefinite cone,
$\mathcal{Y} = \{Y: Y_{00}=1, \mathcal{G}(Y)=0, 0\le Y\le1\}$ is a box with fixed
entries (so $\Pi_{\mathcal{Y}}$ is entrywise clipping), $\beta>0$ is the penalty
and $\gamma\in(0,\frac{1+\sqrt5}{2})$ the dual step.

\paragraph{Published bounds.} The values of \eqref{eq:dnn} reported in the
literature are obtained by running a first-order method to a tolerance, typically
$10^{-5}$ in relative residuals, and evaluating a dual function at the final
iterate, rounded up for integer data \cite{owx,graham,kkt,bbcpop}. Graham et
al.\ \cite{graham} strengthen the subproblems with the redundant constraint
$\tr R = m+1$ and with the constraints of the polyhedral set, which improves the
dual function and the convergence but not the value of the relaxation. For a
branch-and-bound algorithm, in which a bound is needed after a few hundred
iterations at most and an invalid bound invalidates the proof, we need more: a
number that is below the optimum whatever the state of the iteration.

\section{A certified bound}\label{sec:cert0}

\subsection{A bound from any multiplier}\label{sec:cert}

ADMM iterates are neither primal feasible nor dual optimal, and the primal value
$\inner{L}{Y}$ is not a bound: stopped early, it can lie above the optimum.
Branch and bound needs a number that is provably below the subproblem optimum after
any number of iterations. We obtain one from the multiplier.

\begin{lemma}\label{lem:trace}
Every lifted permutation $Y = yy^T$ satisfies $Y = VRV^T$ with
$R = V^TYV\succeq0$ and $\tr(R) = m+1$.
\end{lemma}
\begin{proof}
Since $y$ lies in the range of $V$,
\[
  Y = VV^TYVV^T = VRV^T, \qquad R = V^TYV = (V^Ty)(V^Ty)^T\succeq0,
\]
and $\tr R = \|V^Ty\|^2 = \|y\|^2 = 1 + \sum x_{ij}^2 = m+1$.
\end{proof}

\begin{proposition}\label{prop:cert}
Let $Z\in\Sym^N$ be any symmetric matrix, $\lambda \ge \lmax(V^TZV)$, and let
$\mathcal{P}\supseteq\{yy^T : X\in\Pi\}$. Then the subproblem optimum is at least
\begin{equation}\label{eq:cert}
  s + \min_{Y\in\mathcal{P}} \inner{L+Z}{Y} - (m+1)\lambda .
\end{equation}
\end{proposition}
\begin{proof}
For a lifted permutation $Y = VRV^T$ (Lemma~\ref{lem:trace}), since $R\succeq0$,
\[
  \inner{Z}{Y} = \inner{V^TZV}{R} \le \lmax(V^TZV)\tr R \le (m+1)\lambda .
\]
Hence
\begin{align*}
  \inner{L}{Y} = \inner{L+Z}{Y} - \inner{Z}{Y} &\ge \inner{L+Z}{Y} - (m+1)\lambda\\
  &\ge \min_{\mathcal{P}}\inner{L+Z}{\cdot} - (m+1)\lambda. \qedhere
\end{align*}
\end{proof}

The same bound appears as the dual function of Graham et al.\ \cite[eq.~(3.11),
Thm.~3.4]{graham}. Proposition~\ref{prop:cert} holds for every $Z$; the quality of
the bound depends on $Z$ being near a dual optimum of \eqref{eq:dnn}, and on
choosing $Z$ so that $\lambda$ can be small. Two choices of $\mathcal{P}$ and one of
$Z$ make it both cheap and tight.

\paragraph{The multiplier.} The approach of \cite{owx} projects $V^TZV$ onto the
negative semidefinite cone, which costs an extra eigendecomposition per check.
Instead we use the multiplier implied by the $R$-step. Let $W = V^T(Y+Z/\beta)V$ be
the matrix projected in \eqref{eq:rstep}, and
\[
  Z' = Z + \beta(Y - VRV^T), \qquad\text{so that}\qquad V^TZ'V = \beta(W - \Pspd(W)) \preceq 0 .
\]
The last matrix is $\beta$ times the negative part of $W$, negative semidefinite by
construction. This holds even when the eigenvectors used for $\Pspd$ are
inaccurate: if $\Pspd(W)$ is formed as $W$ minus the computed negative part
$\sum_{w_k<0} w_k v_kv_k^T$, the difference is exactly that sum. In floating
point, $\lmax(V^TZ'V)$ is a tiny positive number, which we bound below.

\paragraph{The polyhedral part.} With $\mathcal{P}$ the box $\mathcal{Y}$ itself,
$\min_{\mathcal{P}}\inner{M}{Y}$ for $M = L+Z'$ is $M_{00}$ plus the sum of the
negative entries of $M$ outside the gangster set, the simplest choice. A much
tighter, still cheap choice uses that every permutation is an assignment in two
levels. Write $M_{(ij),\cdot}$ for row $(ij)$ of $M$. For a permutation $\sigma$,
\begin{equation}\label{eq:gl}
\inner{M}{yy^T} = M_{00} + \sum_{i}\Big( 2M_{0,(i\sigma_i)} + M_{(i\sigma_i),(i\sigma_i)}
 + \sum_{k\ne i} M_{(i\sigma_i),(k\sigma_k)} \Big).
\end{equation}
The inner sum is a perfect matching between facilities $k\ne i$ and locations
$l\ne\sigma_i$, so it is at least the value $\LAP_{ij}$ of the $(m-1)\times(m-1)$
assignment problem on the block $M_{(ij),(kl)}$, $k\ne i$, $l\ne j$. Hence
\begin{equation}\label{eq:glbound}
\min_{\sigma}\inner{M}{yy^T} \;\ge\; \GL(M) := M_{00} + \LAP\big(c\big),\qquad
c_{ij} = 2M_{0,(ij)} + M_{(ij),(ij)} + \LAP_{ij},
\end{equation}
a Gilmore--Lawler computation \cite{gilmore,lawler} on the lifted matrix. It also
bounds the DNN relaxation itself, not only permutations: for feasible $Y$ of
\eqref{eq:dnn}, column $(ij)$ restricted to $k\ne i$, $l\ne j$ is nonnegative with
all row and column sums $x_{ij} = Y_{0,(ij)}$, so it is $x_{ij}$ times a doubly
stochastic matrix, and assignment duality applies twice. It is never below the box
value and reaches a given value 1.5--6$\times$ sooner in the nug20 tests. Any lower
bound on each $\LAP_{ij}$ may replace it, for example the row-then-column
reduction, which costs $O(m^2)$. The implementation uses the reduction, plus exact
Hungarian solves when the result is within a small gap of the cutoff.

\begin{remark}
Nothing in the certificate depends on the ADMM having converged, on the accuracy
of its eigenvectors, or on the single-precision arithmetic used in the $R$-step
(Section~\ref{sec:impl}). Only the multiplier $Z'$, the Cholesky test below and
\eqref{eq:glbound} enter the certified value, all in double precision.
\end{remark}

\subsection{Floating-point safety}\label{sec:fp}

We assume IEEE binary64 arithmetic with rounding to nearest, an extended type
(\texttt{long double}) with a $p$-bit significand ($p=64$ on x86), and no overflow
or underflow; a non-finite value anywhere yields no certificate. The certified
value is $s + \GL(L+Z') - \mu$, rounded down, where the inputs and the margin $\mu$
cover:
\begin{itemize}
\item[(o)] the subproblem data: $s$ and the reduced linear term $C$, sums of
  products that can cancel, are replaced by lower bounds (exact when all operands
  are multiples of $\frac12$ and the sum of magnitudes is below $2^{p-2}$, as for
  integer data; otherwise the extended-precision sum less $(k+2)\,2^{-p}\sum|t|$
  for $k$ terms $t$, rounded down). Bounds for the lowered subproblem are bounds for
  the exact one, since assignment variables are nonnegative;
\item[(i)] $(m+1)\lambda$ with $\lambda\ge\lmax(V^TZ'V)$, obtained by Rump's
  criterion \cite{rump}: if the floating-point Cholesky factorization of
  $\delta I - S$ succeeds for the computed $S \approx V^TZ'V$, then
  $\lmax(S) \le \delta + c$ with an explicit $c$ depending on the dimension, the
  unit roundoff $u$ and $\tr(\delta I - S)$. To this is added
  $64\,m\,u\,\|Z'\|_F$ for the error of computing $S$ through Householder
  transforms rather than with an exact $V$;
\item[(ii)] the summation and rounding errors of \eqref{eq:glbound}, accumulated in
  extended precision: $8(m+2)u\,\|M\|_1$, plus $(m+2)u\sum_{(ij),(kl)}
  \frac12(|a_{ik}b_{jl}|+|a_{ki}b_{lj}|)$ for forming and scaling $L$.
\end{itemize}
These are standard error bounds with generous constants, not interval arithmetic,
and not every constant is derived here: the allowance $64\,m\,u\|Z'\|_F$ for the
Householder transforms rests on a standard backward-error argument and was checked
empirically (errors at most $2.1\,u\|Z'\|_F$), not proved. At $m=28$ the largest
margin seen was about $10^{-7}$ in objective units, against integer objective
values.

\subsection{Implementation}\label{sec:impl}

The bound is implemented in C++ with LAPACK and BLAS.

\paragraph{Products with $V$.} Forming $V$ explicitly and computing $V^TMV$ and
$VRV^T$ with dense products of order $N\times N$ costs $O(m^6)$. We take for
$V_m$ the first $m-1$ columns of the Householder reflector $H$ with
$He_{m-1} = e/\sqrt m$. Then $V^TMV$ and $VRV^T$ are a reflection along each of the
four tensor modes plus a $2\times2$ rotation coupling the constant coordinate,
$O(m^4)$ per product. Both the $R$-step and the certificate are invariant under
an orthogonal change of basis of the face, so any orthonormal $V$ gives the same
iterates and bound.

\paragraph{Eigensolver.} Each $R$-step needs the projection of a $K\times K$
symmetric matrix ($K = 730$ at $m=28$, $2{,}402$ at $m = 50$) with hundreds of
eigenvalues of each sign, so a partial spectrum does not help directly. Between
iterations the matrix changes little. We track its negative eigenspace by a
Chebyshev-filtered subspace iteration in single precision (BLAS level 3), with a
full single-precision eigendecomposition at intervals (every 20 to 120
iterations). Because $R$ is formed as $W$ minus the computed negative part, the
certificate is unaffected by the accuracy of the subspace. When a check fails the
Cholesky test, or the semidefinite penalty alone decides whether the bound
exceeds a given cutoff, the multiplier is re-projected in double precision before
certifying. A similar design, polynomial-filtered projections with a tracked
subspace in reduced precision, has since been proposed for first-order SDP solvers
on GPUs \cite{kang,sunkang}.

\paragraph{The iteration.} The objective is scaled to $\max|L| = 1$. We use the
Peaceman--Rachford form of the method \cite{graham}, with a second dual update
between the $R$-step and the $Y$-step and $\gamma_1 = \gamma = 1$; the certificate
is read from the $R$-step multiplier before that update. The penalty is
$\beta = \theta\, m/3$ with $\theta$ between 0.2 and 0.3 \todo{state the schedule
actually used in Section~\ref{sec:results}, including the separate penalty of the
cut phase (\texttt{dnn.beta\_sched})}. A certificate is computed every 10
iterations.

\todo{Table: cost per iteration by $m$ on a named machine, and the split between
the $R$-step, the $Y$-step and the certificate. The Nov 5 report's
Table~1 compared the C++ code with the numpy prototype, which is of no interest
here; it is dropped.}

\section{Regional cuts}\label{sec:cuts}

The DNN relaxation is weak in a specific way on the problems that interest us.
On a grid problem such as nug30 or wil50 its optimal $Y$ nearly satisfies every
triangle inequality of the Boolean quadric polytope on its own entries: at the
root of nug30 the largest violation is $0.004$, and adding these inequalities
raised the bounds of two difficult nug30 subproblems by 0.2 and 0.4, against gaps
of 76 and 29. % SOURCE: experiments/dnn_cuts/README.md
The slack of the relaxation is not in triples of individual assignments but in
aggregates of them.

\subsection{Regional moments}

Let $S$ be a set of locations. For a permutation, let
\[
  z_i^S = \sum_{a\in S} x_{ia} \in\{0,1\}
\]
indicate that facility $i$ is assigned to a location in $S$. The first and second
moments of these indicators are linear in the lifted matrix $Y = yy^T$:
\begin{equation}\label{eq:moments}
  p_i^S = \sum_{a\in S} Y_{(ia),(ia)} = z_i^S, \qquad
  W_{ij}^S = \sum_{a,b\in S} Y_{(ia),(jb)} = z_i^Sz_j^S \quad (i\ne j).
\end{equation}
For any $Y$ feasible for \eqref{eq:dnn} we take \eqref{eq:moments} as the
definition of $p^S$ and $W^S$. The matrix
$\begin{pmatrix}1 & (p^S)^T\\ p^S & W^S\end{pmatrix}$, with $W^S_{ii} = p^S_i$, is
a linear compression of $Y$ and so is positive semidefinite for every feasible
$Y$: requiring this adds nothing. What the relaxation does not know is that the
$z_i^S$ are binary.

\begin{proposition}\label{prop:tri}
Let $S$ be a set of locations and $i,j,k$ distinct facilities. Every lifted
permutation satisfies
\begin{align}
  W_{ij}^S + W_{ik}^S - W_{jk}^S - p_i^S &\le 0, \label{eq:t1}\\
  p_i^S + p_j^S + p_k^S - W_{ij}^S - W_{ik}^S - W_{jk}^S &\le 1, \label{eq:t2}
\end{align}
and the two inequalities obtained from \eqref{eq:t1} by exchanging $i$ with $j$
or with $k$.
\end{proposition}
\begin{proof}
With $z = (z_i^S, z_j^S, z_k^S)\in\{0,1\}^3$, the left side of \eqref{eq:t1} is
$z_iz_j + z_iz_k - z_jz_k - z_i = -z_i(1-z_j)(1-z_k) - (1-z_i)z_jz_k \le 0$, and
\eqref{eq:t2} is $(1-z_i)(1-z_j)(1-z_k) + z_iz_jz_k \ge 0$ expanded.
\end{proof}

These are the triangle inequalities of the Boolean quadric polytope
\cite{padberg} applied to $z^S$. The same construction applies with the roles of
facilities and locations exchanged ($S$ a set of facilities, and $i,j,k$
locations), and at a node of a branch-and-bound tree with $S$ restricted to the
free locations. We call \eqref{eq:t1}--\eqref{eq:t2} \emph{regional} triangle
inequalities and $S$ a \emph{region}. The case $|S| = 1$ gives the triangle
inequalities on the entries of $Y$.

\subsection{The inequalities are not implied}

\begin{proposition}\label{prop:strict}
There is a QAP with nonnegative data for which the DNN relaxation \eqref{eq:dnn}
has value at most 9, the optimum is 10, and \eqref{eq:dnn} together with a single
inequality \eqref{eq:t2} has value 10.
\end{proposition}
\begin{proof}
Take $n = 6$, split the facilities into two groups of three and the locations
into two groups of three, and let $A = B$ be the adjacency matrix of $K_{3,3}$:
entries 1 across the groups and 0 within. If $t$ facilities of the first group
occupy the first group of locations, the objective is
$2(t^2 + (3-t)^2) = 9 + 4(t-\tfrac32)^2$, so the optimum is 10. Let $S$ be the
first group of locations and define $Q\in\Sym^6$ with unit diagonal, $-\frac12$
between facilities of the same group and 0 between groups, $p_i = \frac12$ and
$W = \frac14(ee^T + Q)$. Set $Y_{00} = 1$,
$Y_{0,(ia)} = Y_{(ia),(ia)} = \frac16$, zero on the gangster entries, and for
$i\ne j$, $a\ne b$,
\[
  Y_{(ia),(jb)} = \begin{cases} W_{ij}/6, & a, b \text{ in the same group of locations},\\
  (\frac12 - W_{ij})/9, & \text{otherwise}. \end{cases}
\]
This $Y$ is feasible for \eqref{eq:dnn} (its eigenvalues are $2$, $\frac14$ eight
times, $\frac3{16}$ sixteen times and $0$ twelve times), satisfies every triangle
inequality on its own entries with slack at least $\frac1{12}$, and has objective
value 9. For the three facilities of the first group it violates \eqref{eq:t2} by
$\frac18$. Finally, with $g = 1 - z_1 - z_2 - z_3 + z_1z_2 + z_1z_3 + z_2z_3$
the slack of that inequality, the objective equals $10 + 8g$ on permutations, so
with multiplier 8 the Lagrangian is constant and equal to 10.
\todo{The feasibility and the spectrum were verified in exact rational
arithmetic (\texttt{experiments/wil50\_feasibility/stronger\_bounds}); recheck
the statement line by line before submission, and decide whether the verification
script is cited or an explicit factorization given.}
\end{proof}

The example shows the mechanism: the relaxation can behave as if $\frac32$
facilities of a group were in $S$, with zero variance. The regional inequalities
forbid a small variance around a fractional count.

\subsection{Choosing the regions}\label{sec:regions}

Proposition~\ref{prop:tri} holds for any $S$. Which regions are useful depends on
the distance matrix.

\paragraph{Grid problems.} For the nug, sko, tho and wil problems of QAPLIB, $B$
is the matrix of Manhattan distances between the cells of an $r\times c$ grid.
Let $\mathcal{S}$ be the $r-1$ row prefixes $\{a : \mathrm{row}(a) < t\}$ and the
$c-1$ column prefixes $\{a: \mathrm{col}(a) < t\}$. Then
\begin{equation}\label{eq:decomp}
  b_{ab} = \sum_{S\in\mathcal{S}} \big|\mathbf{1}_{a\in S} - \mathbf{1}_{b\in S}\big|,
  \qquad\text{so}\qquad
  \sum_{i,j} a_{ij}\,b_{\pi(i)\pi(j)} = \sum_{S\in\mathcal{S}}\sum_{i,j} a_{ij}\big(z_i^S + z_j^S - 2z_i^Sz_j^S\big).
\end{equation}
The objective is a sum over the $r + c - 2$ regions of the flow cut by the set of
facilities placed in the region: a stack of graph-partition problems with
prescribed part sizes, coupled by the requirement that one assignment produce all
the partitions. The regional moments are exactly the variables of these partition
problems, and \eqref{eq:t1}--\eqref{eq:t2} are the triangle inequalities that
strengthen SDP relaxations of graph partition
\cite{dmswz}. For wil50 ($5\times10$) there are 13 regions and
$13\cdot4\binom{50}{3}\approx10^6$ inequalities at the root, few enough to
separate by enumeration.

\paragraph{Clustered distances.} In Taillard's problems tai$n$b \cite{taillard}
the distances are rounded Euclidean distances between points in the plane that lie
in a few tight, well separated clusters, so that 92--98\% of the total distance is
between clusters; on tai20b, tai25b, tai30b and tai50b the clusters are moreover
nearly collinear. % SOURCE: experiments/bound_lab/regions/README.md section 1
A metric on a line decomposes into the cut semimetrics of the prefixes of the
sorted order exactly as \eqref{eq:decomp}, and a planar Euclidean metric
approximately into those of half-planes. We use as regions the prefixes of the
locations sorted along the first principal axis of a two-dimensional scaling of
$B$, the clusters of an average-linkage dendrogram of $B$ and their unions, and
the half-planes in a few further directions that fall between clusters (86
regions for tai60b). \todo{State the family precisely (\texttt{gen\_regions.py},
family \texttt{comb}) and give the number of regions per problem.} Regions on the
facility side, from the flow matrix, and random regions gave no improvement on
these problems.

\paragraph{Related work.} Mittelmann and Peng \cite{mp} use the structure of
Manhattan and Hamming distance matrices through a different fact, the
positive semidefiniteness of $\alpha ee^T - XBX^T$, to obtain small SDP
relaxations; that constraint is implied by \eqref{eq:dnn}, and their bounds are
below the DNN bound (wil50: 46,526). Box inequalities for the QAP polytope
\cite{jk} are inequalities of the form $(q - \beta)(q-\beta-1)\ge0$ on integer
counts $q$ of assignments. The regional triangle inequalities with three
facilities are related to them, and we tested the larger members of that family
(Section~\ref{sec:other}). De Meijer et al.\ \cite{dmswz} add triangle
inequalities to DNN relaxations of the graph partition problem by a cutting-plane
ADMM; the compression of \eqref{eq:dnn} to one region is such a relaxation of a
bisection problem. We have not found the regional inequalities, with regions
taken from a decomposition of the distance matrix, used for the QAP.
\todo{Literature check to be repeated before submission
(\texttt{experiments/bound\_lab/literature}; Gueye and Michelon 2015 was read in
abstract only).}

\subsection{The certificate with cuts}\label{sec:certcuts}

Write the inequalities in use as $\inner{G_t}{Y}\le h_t$, $t = 1,\dots,T$, with
$G_t\in\Sym^N$ sparse and $h_t\in\{0,1\}$.

\begin{proposition}\label{prop:certcuts}
For any $Z\in\Sym^N$, $\lambda\ge\lmax(V^TZV)$ and $\mu\in\R^T$, $\mu\ge0$, the
subproblem optimum is at least
\begin{equation}\label{eq:certcuts}
  s + \GL\Big(L + \sum_t\mu_tG_t + Z\Big) - \sum_t \mu_th_t - (m+1)\lambda .
\end{equation}
\end{proposition}
\begin{proof}
For a lifted permutation $Y$, $\inner{L}{Y}\ge\inner{L+\sum_t\mu_tG_t}{Y} -
\sum_t\mu_th_t$ by Proposition~\ref{prop:tri} and $\mu\ge0$. Apply
Proposition~\ref{prop:cert} and \eqref{eq:glbound} to the lifted cost
$L + \sum_t\mu_tG_t$.
\end{proof}

So a bound with cuts is the certified bound of Section~\ref{sec:cert0} for a
modified cost matrix, less a constant. The floating-point margins of
Section~\ref{sec:fp} are extended by the magnitudes of the added terms, with
$\sum_t\mu_th_t$ rounded up.

\subsection{Cuts inside the ADMM}\label{sec:cpadmm}

Our first implementation used \eqref{eq:certcuts} directly: solve \eqref{eq:dnn}
with cost $L + \sum\mu_tG_t$, separate the most violated inequalities at the final
$Y$, update $\mu$ by a subgradient step, and repeat. This needs many rounds of
several hundred iterations each. It is far more effective to treat the
inequalities as constraints of the polyhedral step, as in the cutting-plane ADMM
of \cite{dmswz}. The $Y$-step \eqref{eq:ystep} becomes the projection of
$X = VRV^T - (L+Z)/\beta$ onto
\[
  \mathcal{Y}\cap\{Y : \inner{G_t}{Y}\le h_t,\ t = 1,\dots,T\},
\]
computed approximately by Dykstra's algorithm \cite{dykstra}, cycling over
$\mathcal{Y}$ and the half-spaces. The correction that Dykstra's algorithm keeps
for half-space $t$ is a multiple $\nu_tG_t$ of its normal with $\nu_t\ge0$, and
the step for it is Hildreth's,
\[
  \nu_t' = \max\big(0,\ \nu_t + (\inner{G_t}{Y} - h_t)/\|G_t\|^2\big),\qquad
  Y \leftarrow Y - (\nu_t'-\nu_t)\,G_t .
\]
Since $\mathcal{Y}$ is visited first in each sweep, its step is
$Y = \Pi_{\mathcal{Y}}(X - \sum_t\nu_tG_t)$ regardless of its own previous
correction, so the state of the projection is the vector $\nu$ alone, which is
carried from one ADMM iteration to the next. We use one sweep per iteration.

\paragraph{Cost.} An inequality \eqref{eq:t1}--\eqref{eq:t2} is linear in the
moments $(p^S, W^S)$ of its region, and a step along $G_t$ changes the moments of
another region $S'$ by multiples that depend only on $|S\cap S'|$. The half-space
steps therefore run on the moment matrices of the regions, about 80 floating-point
operations per inequality, and $Y$ is touched once per sweep, to form
$X - \sum\nu_tG_t$, clip it, and accumulate the moments of the result. With a
pool of 13,000 inequalities at a wil50 subproblem with $m = 47$ the projection
adds 4\% to 13\% to the cost of an iteration, depending on how often the
certificate is evaluated. % SOURCE: experiments/wil50_feasibility/roadmap/admm/cpadmm/README.md

\paragraph{The certificate.} The optimality conditions of the projection give
multipliers $\mu = \beta\nu$ for the inequalities, in the units of the scaled
cost (times $\max|L|$ in the original units), and $\mu\ge0$ by construction. At
each check the bound \eqref{eq:certcuts} is evaluated with this $\mu$ and the
multiplier $Z'$ of Section~\ref{sec:cert}. As before, nothing in the certified
value depends on the projection having converged.

\paragraph{Separation.} Every 100 iterations, the 1,000 most violated regional
triangle inequalities at the current (projected) $Y$ are added to the pool with
$\nu_t = 0$; separation is by enumeration over the regions and triples.
\todo{Describe the ranking, and whether inequalities are ever dropped in the
settings used for the results.}

\paragraph{Effect on convergence.} On a wil50 subproblem at depth 3, with the
same pool of 13,000 inequalities, the cutting-plane form reaches a given
certified value in 4.7 times fewer iterations than the subgradient form (1,040
against 4,860) and continues to a higher value; on two further subproblems the
ratios are 5.7 and 4.8. % SOURCE: cpadmm/README.md, "Short answer"

\subsection{Other inequalities}\label{sec:other}

The triangle inequalities are the smallest members of larger families on the
same moments: the inequalities $(q-\beta)(q-\beta-1)\ge0$ for the count
$q = \sum_{i\in T}z_i^S$ of a set $T$ of facilities (for $|T|=5$ these include the
pentagonal inequalities of \cite{padberg}; for general $T$, the box inequalities
of \cite{jk}), and all valid inequalities on seven indicators, separated by
projection onto the convex hull of the feasible $0/1$ configurations. We
implemented these and found them not worth their cost after the triangle
inequalities on wil50 and on tho40. We also tested inequalities valid only for
permutations that no pairwise exchange improves; they are valid for an optimal
solution and their exact strength is considerable, but the DNN relaxation
realized about a tenth of it. \todo{One sentence each with the measured gains
(\texttt{bound\_lab/tho40\_nodes}, \texttt{bound\_lab/swapcuts},
\texttt{wil50\_feasibility/new\_cuts}), or drop the subsection.}

\section{Use in branch and bound}\label{sec:bnb}

A bound for branch and bound is judged by more than its value at the root. At a
node we want to know, as cheaply as possible, whether the bound exceeds the
incumbent value less a tolerance (the \emph{cutoff}), in which case the node is
fathomed; and if it does not, we want what was computed to be of use to the
node's children. We describe the properties of $\DNNp$ that serve these ends.
How they are combined with a cheaper bound, with branching rules and with a
parallel search is described in \cite{companion}.

\paragraph{Stopping at the cutoff.} Because the certified value is a bound at
every check, the iteration stops at the first check at which it exceeds the
cutoff. Most fathomed nodes need a small fraction of the iterations that
convergence would. No tolerance on residuals is involved.

\paragraph{Giving up.} Conversely, a node whose certified value, extrapolated
generously from its recent gains, stays below the cutoff can be abandoned early
and branched on. This affects the cost of the search but not its validity.

\paragraph{Warm starts.} A child subproblem fixes one more assignment
$x_{ij} = 1$. Its ADMM is started from the parent's final multiplier $Z$ with the
row and column of the new assignment removed from the lifted indices, and from
the parent's final $Y$ conditioned on the assignment: the parent's row $(ij)$,
divided by its diagonal entry, becomes the child's row 0, and the remaining
entries are rescaled so that the diagonal of $Y$ equals that row. The certificate
does not depend on the start. \todo{Measured effect with the current code; the
Nov 5 report had 40\% fewer iterations from $Z$ alone on nug30 and a factor 1.3
from adding $Y$ on 16 wil50 children.} The parent's state is stored with the
node in single precision (the upper triangles of $Y$ and $Z$), so that a child
can be bounded in another process or on another machine.

\paragraph{Inheriting cuts.} Proposition~\ref{prop:tri} holds for any region and
any three free facilities, so an inequality of the parent's pool is valid for a
child as long as none of its three facilities is the one just fixed; the region is
restricted to the child's free locations, and the right-hand side is unchanged.
A child therefore starts with the parent's active inequalities and their
multipliers, in addition to $Y$ and $Z$, and a parent that is not fathomed still
equips its children. \todo{Measured effect of inheritance on a fixed set of
children, for example tho40 or wil50 depth-4 nodes.}

\paragraph{Two phases.} At a node we first run the ADMM without new cuts (with
the inherited pool, if any), and enter the cutting-plane phase only if the node
is not fathomed and its bound is within a given fraction of the cutoff.
\todo{Decide how much of the policy (zone, hand-over rule, phase penalty) is
described here and how much in the companion paper. Proposal: the two-phase
structure and the zone only.}

\paragraph{An illustration.} \todo{One table on problems solved by others, to
show the effect per node without announcing new solutions: for example the
fraction of depth-$d$ subproblems of nug30 and sko42 fathomed by DNN and by
$\DNNp$ within a fixed iteration budget; and tai35b and tai40b, whose trees
under $\DNNp$ have 922 and 733 nodes (\texttt{bound\_lab/tai\_solves}) against
$2.6\times10^6$ and $2.8\times10^5$ in \cite{fujii}. To be agreed.}

\section{Computational results}\label{sec:results}

\todo{Machine, compiler, BLAS, threads; the stopping rule for root bounds (a
fixed iteration limit, or a stall rule on the certified value); the settings
file for each row. All root bounds below were computed in different sessions with
different limits and several were still rising when stopped; they are to be
recomputed under one rule (\texttt{docs/dnn\_plus/root\_runs.md}).}

\subsection{The certified DNN bound on small and medium problems}

In Table~\ref{tab:bounds} we give, for the problems of \cite[Table~1]{ab01}, the
optimum value and the lower bounds GLB, PB, QPB and SDPB1 reported there, together
with the certified DNN bound of Section~\ref{sec:cert0}, without cuts, at the
root. GLB, PB and SDPB1 (the cheapest SDP relaxation of \cite{zkrw}) are taken
from \cite{zkrw}, and QPB is the bound computed in \cite{ab01} with one step of
its dual improvement. The DNN bound was computed with up to 4,000 ADMM iterations.
All bounds are rounded up to the next integer, since the data for these problems
is integral.

\begin{table}[p]\centering
\caption{Bounds for QAPLIB problems. GLB, PB, QPB and SDPB1 from
\cite[Table~1]{ab01}; DNN: the certified bound of Section~\ref{sec:cert0} at the
root, after at most 4,000 ADMM iterations; --: held out (see text).}\label{tab:bounds}
\small
\begin{tabular}{lrrrrrr}
\toprule
Problem & OPT & GLB & PB & QPB & SDPB1 & DNN\\
\midrule
esc16a & 68 & 38 & 47 & 55 & 47 & 64\\
esc16b & 292 & 220 & 250 & 250 & 250 & 290\\
esc16c & 160 & 83 & 95 & 95 & 95 & 154\\
esc16d & 16 & 3 & $-$19 & $-$19 & $-$19 & 13\\
esc16e & 28 & 12 & 6 & 6 & 6 & 27\\
esc16g & 26 & 12 & 9 & 9 & 9 & 25\\
esc16h & 996 & 625 & 708 & 708 & 708 & 977\\
esc16i & 14 & 0 & $-$25 & $-$25 & $-$25 & 12\\
esc16j & 8 & 1 & $-$6 & $-$6 & $-$6 & 8\\
\midrule
had12 & 1652 & 1536 & 1573 & 1592 & 1604 & 1652\\
had14 & 2724 & 2492 & 2609 & 2630 & 2651 & 2724\\
had16 & 3720 & 3358 & 3560 & 3595 & 3612 & 3720\\
had18 & 5358 & 4776 & 5104 & 5143 & 5174 & 5358\\
had20 & 6922 & 6166 & 6625 & 6677 & 6713 & 6922\\
\midrule
kra30a & 88900 & 68360 & 63717 & 68572 & 69736 & --\\
kra30b & 91420 & 69065 & 63818 & 69021 & 70324 & --\\
\midrule
nug12 & 578 & 493 & 472 & 482 & 486 & 568\\
nug14 & 1014 & 852 & 871 & 891 & 903 & 1011\\
nug15 & 1150 & 963 & 973 & 996 & 1009 & 1141\\
nug16a & 1610 & 1314 & 1403 & 1448 & 1461 & 1600\\
nug16b & 1240 & 1022 & 1046 & 1071 & 1082 & 1219\\
nug17 & 1732 & 1388 & 1487 & 1529 & 1548 & 1708\\
nug18 & 1930 & 1554 & 1663 & 1705 & 1723 & 1894\\
nug20 & 2570 & 2057 & 2196 & 2254 & 2281 & 2507\\
nug21 & 2438 & 1833 & 1979 & 2055 & 2090 & 2382\\
nug22 & 3596 & 2483 & 2966 & 3080 & 3140 & 3529\\
nug24 & 3488 & 2676 & 2960 & 3028 & 3068 & 3402\\
nug25 & 3744 & 2869 & 3190 & 3272 & 3305 & 3626\\
nug30 & 6124 & 4539 & 5266 & 5365 & 5413 & 5950\\
\midrule
rou12 & 235528 & 202272 & 200024 & 206102 & 208685 & 235528\\
rou15 & 354210 & 298548 & 296705 & 303777 & 306833 & 350217\\
rou20 & 725520 & 599948 & 597045 & 607822 & 615549 & 695181\\
\midrule
scr12 & 31410 & 27858 & 4727 & 8585 & 11117 & 31410\\
scr15 & 51140 & 44737 & 10355 & 12479 & 17046 & 51140\\
scr20 & 110030 & 86766 & 16113 & 23960 & 28535 & 106802\\
\midrule
tai12a & 224416 & 195918 & 193124 & 199597 & 203595 & 224416\\
tai15a & 388214 & 327501 & 325019 & 330310 & 333437 & 377101\\
tai17a & 491812 & 412722 & 408910 & 416033 & 419619 & 476525\\
tai20a & 703482 & 580674 & 575831 & 585139 & 591994 & 671675\\
tai25a & 1167256 & 962417 & 956657 & 983456 & 974004 & 1096657\\
tai30a & 1818146 & 1504688 & 1500407 & 1518059 & 1529135 & 1706870\\
\midrule
tho30 & 149936 & 90578 & 119254 & 124684 & 125972 & 143576\\
\bottomrule
\end{tabular}
\end{table}

Over the 40 problems other than kra30a and kra30b the DNN bound closes on
average 89\% of the gap between QPB and the optimum, and at least 62\% on every
problem. It equals the optimum on ten of them. \todo{kra30a and kra30b: the Nov 5
table held these out because the ADMM did not converge with the default penalty.
\texttt{experiments/kra30\_solve} later quotes root bounds of 86,837 and 87,858;
confirm and fill in. Add a DNN$^+$ column for the grid problems (nug, tho30) and
decide whether this table keeps PB and SDPB1.}

\subsection{Large problems}

Table~\ref{tab:large} compares $\DNNp$ at the root with the best published lower
bounds for the larger grid and tai$n$b problems of QAPLIB. The published values
are all bounds from the DNN relaxation \eqref{eq:dnn}: rPRSM is the restricted
Peaceman--Rachford method of \cite{graham}, NB the Newton-bracketing method of
\cite{kkt}, and BBCPOP the bisection and projection method of \cite{bbcpop}.

\begin{table}[t]\centering
\caption{Root lower bounds for large QAPLIB problems. BKV: best known value.
Published DNN bounds from \cite{graham} (rPRSM), \cite{kkt} (NB) and \cite{bbcpop}
(BBCPOP); --: not reported. DNN and $\DNNp$: the certified bounds of this paper
without and with regional cuts. Gap: $(\mathrm{BKV}-\DNNp)/\mathrm{BKV}$.}\label{tab:large}
\small\setlength{\tabcolsep}{4pt}
\begin{tabular}{lrrrrrrr}
\toprule
Problem & BKV & rPRSM & NB & BBCPOP & DNN & $\DNNp$ & Gap\\
\midrule
sko42 & 15,812 & 15,336 & 15,334 & -- & \todo{run} & \todo{run} & \\
sko49 & 23,386 & 22,654 & 22,653 & 22,650 & \todo{run} & \todo{run} & \\
tho40 & 240,516 & 226,522 & 226,516 & 226,490 & \todo{run} & 230,453 & 4.18\%\\
wil50 & 48,816 & 48,126 & 48,123 & 48,121 & \todo{run} & \todo{run} & \\
\midrule
tai25b & 344,355,646 & -- & -- & \todo{look up} & 337,285,170 & 344,084,304 & 0.08\%\\
tai30b & 637,117,113 & -- & -- & 598,629,428 & 597,633,906 & 636,993,730 & 0.02\%\\
tai35b & 283,315,445 & -- & -- & \todo{look up} & 269,414,422 & 279,325,031 & 1.41\%\\
tai40b & 637,250,948 & -- & -- & \todo{look up} & 608,514,446 & 627,372,144 & 1.55\%\\
tai50b & 458,821,517 & -- & 431,245,612 & 431,090,745 & 430,147,161 & 453,552,676 & 1.15\%\\
tai60b & 608,215,054 & -- & 592,829,061 & -- & \todo{run} & 605,315,696 & 0.48\%\\
\bottomrule
\end{tabular}
\end{table}
% SOURCE: published columns: experiments/bound_lab/literature/README.md (Q1 table).
% tho40 230,452.99: bound_lab/cpfast (logs/R_wc288_long.log), 300 + 2,500 iterations.
% tai25b-40b: bound_lab/regions/README.md section 6, 2,300 total iterations; tai30b DNN+ at 6,300.
% tai35b, tai40b: the 6,000-iteration roots of bound_lab/tai_solves are 0.89% and 0.18% under the optimum,
%   higher than the values above; use those once the exact figures are read from the logs.
% tai50b: bound_lab/tai50b_size (470 + 6,000 iterations, 4,341 CPU s); plain 430,147,161 at 1,500 iterations.
% tai60b: docs/next_runs.md (605,315,696.64), computed on branch claude/cool-babbage-ilnogj.

\todo{Discussion, once the table is complete. Points the present numbers support:
on the tai$n$b problems regional cuts close 66\% to 99.7\% of the gap left by the
DNN relaxation; on tho40 they close about 28\% of it (from 226,522 to 230,453
against 240,516), and the bound was still rising. Points not yet supported: sko42,
sko49 and wil50 have no root $\DNNp$ value; the flows of wil50 are 90\% dense and
of sko49 69\%, against 40\% for tho40, and triangle inequalities on partition
relaxations are known to gain less on dense graphs \cite{dmswz}, so the gain there
may be small and should be reported as it is.}

\subsection{Cost}

\todo{Table: iterations and CPU time for each row of Table~\ref{tab:large}, with
the time per iteration, on one machine. For orientation, the published times for
the DNN relaxation of tho40 are 2,323 s (rPRSM, 6,200 iterations), about 5,000 s
(NB), and 18,000 s (BBCPOP) on their authors' machines; the tai50b value above
took 4,341 CPU seconds. A figure of certified bound against iterations for one
grid problem and one tai$n$b problem, with and without cuts, would show both the
any-time property and the effect of the cuts.}

\section{Conclusion}

\todo{To be written when the results are complete. Themes: the DNN relaxation's
slack on structured QAPs lies in aggregates that a decomposition of the distance
matrix identifies; imposing integrality of those aggregates through triangle
inequalities is cheap inside a first-order method; certification at every
iteration makes the bound usable in branch and bound. Open: regions for problems
without metric structure (tai$n$a), denser flows, a verified evaluation of the
certificate in interval arithmetic, a GPU implementation.}

\appendix
\section{Validation}\label{sec:validation}

A branch-and-bound proof is only as good as its bounds, so the DNN bound was
checked in several independent ways.
\begin{itemize}
\item \textbf{Exhaustive checks.} 1,512 certified bounds on random instances and
  subproblems with $m\le8$ (several families, including linear terms and ties),
  each compared with the optimum found by enumeration. No bound exceeds the
  optimum, and no fathom is ever claimed with the cutoff at the optimum.
\item \textbf{Inside the search.} The solver's brute-force test (4,536 runs on
  random instances with $n\le 8$, eight parameter sets, three upper-bound modes)
  passes with the DNN bound applied at every depth, with and without the give-up
  rule.
\item \textbf{Against the prototype.} In prototype mode the C++ bound reproduces the
  numpy prototype's certified values at 44 checkpoints to a relative $8.6\cdot10^{-8}$.
  On nug20 nodes whose subproblems were solved exactly, the bounds are below the
  optima and fathom the same nodes.
\item \textbf{Reproducibility.} Results are bitwise identical across runs on one
  machine. Valgrind reports no memory errors on a debug build. Serial, resumed and
  parallel runs give identical node counts.
\item \textbf{Independent mathematical review.} A separate AI agent, given the code
  and this claim but not our arguments, re-derived every step: the face basis
  actually used (Householder reflector and rotation) has the required range, so
  Lemma~\ref{lem:trace} holds for it; Proposition~\ref{prop:cert} holds for any
  sign of $\lmax$; the Gilmore--Lawler evaluation counts every quadratic term
  once and is valid for any dual potentials, so neither the inexact inner LAPs
  nor the outer Hungarian solve need to be exact; the lifting matches the
  solver's objective convention without symmetry assumptions; a NaN or infinite
  multiplier cannot produce a fathom; and the search uses the DNN value only to
  fathom. It found no proof-breaking issue.
\item \textbf{Independent code review and stress tests.} A second agent reviewed
  the implementation and ran 2,592 bounds on brute-forced instances with
  $m\le7$ (asymmetric and negative data, ties, scales up to $10^6$), 540
  warm-started children, an independent long-double check of the face
  transforms (error at most $2.1\,u\|Z\|_F$ against the $64\,m\,u\|Z\|_F$
  allowed), 21,000 random assignment problems against the Hungarian solver,
  and determinism tests, under address, undefined-behavior and leak sanitizers.
  No bound exceeded an optimum and nothing crashed. Its minor findings (and the
  mathematical review's) were fixed: a loop that could skip the last Cholesky
  attempt, widened error constants, an explicit finiteness test on the
  certified value, and skipping the bound while there is no incumbent.
\item \textbf{Adversarial tests.} A third agent tried to break the bound with
  extreme data: magnitudes up to $10^{9}$ and beyond, products above $2^{53}$,
  signed, non-integer and non-symmetric data, rank-one, constant and all-zero
  matrices, subproblems with one to three free facilities, and parameter
  fuzzing; about 15,000 bounds and 8,600 solver runs in all, with an
  instrumented build checking every DNN call inside the search against the
  exact subproblem optimum. On ordinary data no bound was invalid, and the
  rounding part of the margin was never needed (the unmargined value minus the
  semidefinite penalty stayed below the optimum by at least 1.6\% of the rounding
  margin). It did find two ways to exceed the optimum, both needing data that is
  neither symmetric nor small. When neither $A$ nor $B$ is symmetric, an entry of
  $L$ can be the difference of two large products whose rounding errors the
  margin, which scaled with $\sum|L|$, did not cover; the margin now uses the sum
  of the products' magnitudes. And the shift and linear term of a subproblem, sums
  of products that can cancel, were rounded in double before any bound saw them.
  Integer data whose products exceed $2^{53}$ is no longer given an integer
  fathoming tolerance. It also found
  two older bugs outside the DNN bound: a finite sentinel ($10^8$) in the
  assignment solver, which crashed or gave invalid QPB bounds when costs exceeded
  it (fixed), and wrong QPB results on data with negative entries in $A$ or $B$
  (open; QAPLIB data are nonnegative). None of this affects nonnegative integer
  data whose objective values stay below $2^{53}$, where every product and sum is
  exact. That includes all of QAPLIB, even the bur26 family, in which neither
  matrix is symmetric. On the regression instances, node counts are unchanged.
\item \textbf{External review.} An independent review of an earlier draft and the
  code (Astra, an OpenAI Codex agent) re-derived the bound and confirmed that QPB
  values stay paired with their own reduced costs. It showed that accumulating the
  shift and linear terms in extended precision, our first repair, is not enough:
  with $ab = 2^{64}+2^{34}+3$ and $c = -(2^{64}+2^{34})$ the extended-precision
  product rounds by 1 before cancelling, and both the $m=1$ path and a reduced
  subproblem fathomed at a cutoff equal to the optimum. These data now get the
  lower bounds of item (o) in Section~\ref{sec:fp}; both cases, and the
  adversarial ones above, are regression tests (\texttt{tests/dnn\_test.py}, part D)
  that check the bound against optima computed in exact rational arithmetic.
\end{itemize}

\paragraph{Limitations.} The certificate is a theorem in exact arithmetic
(Proposition~\ref{prop:cert} and \eqref{eq:glbound}). In floating point it rests
on the margins of Section~\ref{sec:fp}: standard error bounds with generous
constants under the stated arithmetic model, not interval arithmetic, and not
every constant is derived. The validation above is strong regression evidence, not
a formal verification. For the nug instances the data are small integers, so the
subproblem data are exact, and a false fathom would need a certified value above
the optimum by more than the 0.05 cushion, five to six orders of magnitude above
any margin observed. For non-integer data, where the tolerance is zero, validity
rests on the margins alone; a verified evaluation (directed rounding or interval
arithmetic) of each fathoming certificate would be the safe choice.

\paragraph{Validation of the cuts.} \todo{Rewrite as prose.} The bound with
regional cuts was checked against optima found by enumeration: for the
subgradient form, \todo{count} values; for the cutting-plane form, 12,656
certified values on 77 nodes of random instances with $n = 7, 8, 9$ and of nug12,
none above the optimum (the largest exceeds it by $-1.2\cdot10^{-9}$), and no
fathom with the cutoff at the optimum; for arbitrary regions, 240 node checks and
32 root checks on small tai$n$b and random problems. The assembled cut terms were
compared entry by entry with an independent assembly (relative difference
$10^{-15}$).
% SOURCE: cpadmm/README.md "Validity"; bound_lab/regions/README.md section 3.

\section*{Acknowledgments}
\todo{Nathan to write. The facts, from the repository's record: Claude (an AI
system by Anthropic) wrote the software, ran the experiments and drafted this
report, using independent AI agents for prototype studies, reviews and the
literature search. Astra (an OpenAI Codex agent) proposed imposing consistency
of regional membership indicators on the DNN relaxation, with the
$K_{3,3}$ example of Proposition~\ref{prop:strict} and the first pilot on a wil50
subproblem (memo of October 31, 2026,
\texttt{experiments/wil50\_feasibility/stronger\_bounds/research\_memo\_astra.md}),
and, reviewing an earlier draft, found the cancellation cases of
Appendix~\ref{sec:validation}. The cluster regions for the tai$n$b problems, the
cutting-plane form and its certificate were developed in Claude sessions
(\texttt{experiments/bound\_lab}, \texttt{experiments/wil50\_feasibility/roadmap}).
Journals and arXiv require disclosure of the use of AI tools; this section and
the authors' responsibility statement should meet the target journal's policy.}

{\small
\begin{thebibliography}{99}
% \todo: every entry below marked (verify) was written from memory or from the
% repository's notes and must be checked against the publication.
\bibitem{ab01} K.M. Anstreicher and N.W. Brixius, A new bound for the quadratic
  assignment problem based on convex quadratic programming, \emph{Mathematical
  Programming} 89 (2001) 341--357.
\bibitem{abgl} K. Anstreicher, N. Brixius, J.-P. Goux and J. Linderoth, Solving large
  quadratic assignment problems on computational grids, \emph{Mathematical
  Programming} 91 (2002) 563--588.
\bibitem{ba01} N.W. Brixius and K.M. Anstreicher, Solving quadratic assignment
  problems using convex quadratic programming relaxations, \emph{Optimization
  Methods and Software} 16 (2001) 49--68.
\bibitem{companion} N.W. Brixius and J.-P. Goux, \todo{companion paper: title},
  in preparation.
\bibitem{burkardsurvey} R.E. Burkard, E. \c{C}ela, P.M. Pardalos and L.S. Pitsoulis,
  The quadratic assignment problem, in: D.-Z. Du and P.M. Pardalos (eds.),
  \emph{Handbook of Combinatorial Optimization}, vol.~3, Kluwer (1998) 241--337.
\bibitem{qaplib} R.E. Burkard, S.E. Karisch and F. Rendl, QAPLIB -- a quadratic
  assignment problem library, \emph{Journal of Global Optimization} 10 (1997) 391--403.
\bibitem{datenagi} K. Date and R. Nagi, Level 2 reformulation linearization
  technique-based parallel algorithms for solving large quadratic assignment
  problems on graphics processing unit clusters, \emph{INFORMS Journal on
  Computing} 31 (2019) 771--789. (verify)
\bibitem{dmswz} F. de Meijer, R. Sotirov, A. Wiegele and S. Zhao, Partitioning
  through projections: strong SDP bounds for large graph partition problems,
  \emph{Computers \& Operations Research} 151 (2023) 106088. (verify)
\bibitem{dykstra} R.L. Dykstra, An algorithm for restricted least squares
  regression, \emph{Journal of the American Statistical Association} 78 (1983)
  837--842.
\bibitem{fujii} K. Fujii, N. Ito, S. Kim, M. Kojima, Y. Shinano and K.-C. Toh, Solving
  challenging large scale QAPs, arXiv:2101.09629 (2021).
\bibitem{gilmore} P.C. Gilmore, Optimal and suboptimal algorithms for the quadratic
  assignment problem, \emph{SIAM Journal on Applied Mathematics} 10 (1962) 305--313.
\bibitem{graham} N. Graham, H. Hu, J. Im, X. Li and H. Wolkowicz, A restricted dual
  Peaceman--Rachford splitting method for a strengthened DNN relaxation for QAP,
  \emph{INFORMS Journal on Computing} 34 (2022) 2125--2143.
\bibitem{hrw} S.W. Hadley, F. Rendl and H. Wolkowicz, A new lower bound via projection
  for the quadratic assignment problem, \emph{Mathematics of Operations Research} 17
  (1992) 727--739.
\bibitem{bbcpop} N. Ito, S. Kim, M. Kojima, A. Takeda and K.-C. Toh, Algorithm 996:
  BBCPOP: a sparse doubly nonnegative relaxation of polynomial optimization
  problems with binary, box, and complementarity constraints, \emph{ACM
  Transactions on Mathematical Software} 45 (2019) 34. (verify)
\bibitem{jk} M. J\"unger and V. Kaibel, Box-inequalities for quadratic assignment
  polytopes, \emph{Mathematical Programming} 91 (2001) 175--197.
\bibitem{kang} S. Kang, X. Han, D. Groudiev and H. Yang, \todo{title},
  arXiv:2507.09165 (2025). (verify)
\bibitem{kkt} S. Kim, M. Kojima and K.-C. Toh, A Newton-bracketing method for a
  simple conic optimization problem, \emph{Optimization Methods and Software} 36
  (2021) 371--388. (verify)
\bibitem{lawler} E.L. Lawler, The quadratic assignment problem, \emph{Management
  Science} 9 (1963) 586--599.
\bibitem{loiola} E.M. Loiola, N.M.M. de Abreu, P.O. Boaventura-Netto, P. Hahn and
  T. Querido, A survey for the quadratic assignment problem, \emph{European Journal
  of Operational Research} 176 (2007) 657--690.
\bibitem{mp} H. Mittelmann and J. Peng, Estimating bounds for quadratic assignment
  problems associated with Hamming and Manhattan distance matrices based on
  semidefinite programming, \emph{SIAM Journal on Optimization} 20 (2010)
  3408--3426. (verify)
\bibitem{owx} D.E. Oliveira, H. Wolkowicz and Y. Xu, ADMM for the SDP relaxation of
  the QAP, \emph{Mathematical Programming Computation} 10 (2018) 631--658.
\bibitem{padberg} M. Padberg, The Boolean quadric polytope: some characteristics,
  facets and relatives, \emph{Mathematical Programming} 45 (1989) 139--172.
\bibitem{rump} S.M. Rump, Verification of positive definiteness, \emph{BIT
  Numerical Mathematics} 46 (2006) 433--452.
\bibitem{sunkang} \todo{Sun, Kang and Yang}, arXiv:2609.21102 (2026). (verify)
\bibitem{taillard} E.D. Taillard, Comparison of iterative searches for the
  quadratic assignment problem, \emph{Location Science} 3 (1995) 87--105.
\bibitem{zkrw} Q. Zhao, S.E. Karisch, F. Rendl and H. Wolkowicz, Semidefinite
  programming relaxations for the quadratic assignment problem, \emph{Journal of
  Combinatorial Optimization} 2 (1998) 71--109.
\end{thebibliography}}

\end{document}
